\documentclass[../Probability.tex]{subfiles}

\begin{document}

\section{Estimation}
Given indepedent observations $X_1, \ldots, X_n$ that come from some unknown, true probability distribution $P$, how do we infer $P$?
\begin{enumerate}
    \item Moment Estimation: Estimate the mean and variance of the data.
    \item Finding the full mathematical curve (the distribution) that best fits your dataset, bu tuning parameters. By learning this underlying shape, AI models can sample from it to generate new, realistic synthetic data.
    \item Regression/Classification: Supervised learning, where the observations come in pairs. An input and its corresponding label, denoted as $(X_{i}, Y_{i})$. The goal is to estimate the relationship between them, such as predicting the conditional probability of a label given an input ($P_{Y|X}$) or the expected value ($\mathbb{E}[Y|X]$).
\end{enumerate}

\subsection{Estimator/Estimate}
The primary goal of statistical inference (or learning) is to construct an estimator for $\theta_{*}$.\newline
This $\theta_{*}$ represents a true, unknown parameter associated with a probability distribution $P$.\newline
Examples of this parameter include the mean, the variance, or the parameter of a probabilistic model, $p_{\theta}$.\newline
\textbf{Estimator vs Estimate}:
\begin{itemize}
    \item Estimator: This is defined as a random variable, $\hat{\theta}$, that depends on your data observations ($X_{1},...X_{n}$). It is the overarching mathematical formula or rule, written as $\hat{\theta}=\hat{\theta}(X_{1},...X_{n})$.
    \item Estimate: This is the actual output you get when you plug specific, concrete data points ($x_{1},...,x_{n}$) into your estimator formula. It is the final calculated result, written as $\hat{\theta}(x_{1},...,x_{n})$.
\end{itemize}
The 4 criteria to evaluate if an estimator is good:
\begin{itemize}
    \item \textbf{Unbiasedness}: An estimator $\hat{\theta}(X_{1},...X_{n})$ is unbiased for a true parameter $\theta$ if $\mathbb{E}[\hat{\theta}]=\theta$. The average of many estimates is the true answer.
    \item \textbf{Consistency}: An estimator $\hat{\theta}(X_{1},...X_{n})$ is consistent for a true parameter $\theta$ if $\forall \epsilon > 0,\:\lim\limits_{n\rightarrow\infty}P(|\hat{\theta}-\theta|\le\epsilon)=1$. A single estimate becomes perfectly accurate with enough data.
    \item \textbf{Effectivity}: An unbiased estimator is considered more effective if it has a lower variance than another unbiased estimator, i.e. the confidence is greater.
    \item \textbf{Asymptotic normality}: As $n \to \infty$, distribution of your estimator's guesses will start to perfectly resemble a Normal Distribution (a classic bell curve) centered around the true parameter, by the CLT.
\end{itemize}

\subsubsection{Bias-variance decomposition}
When we create an estimator ($\hat{\theta}$) to guess a true, unknown parameter ($\theta_*$), we need a way to mathematically measure how "good" that guess is.\newline
\textbf{The motivation}: Suppose you have independent data points $X_1, X_2, ..., X_n$ and you want to estimate the true mean.
\begin{enumerate}
    \item $\hat{\theta}_1 = \frac{1}{n} \sum_{i=1}^n X_i$ (The sample mean, averaging $n$ data points).
    \item $\hat{\theta}_2 = X_1$ (Just the first data point).
\end{enumerate}
Both estimators are unbiased, i.e. if you repeat the experiment infinitely many times, both will eventually average out to the true mean.\newline
However, intuitively, $\hat{\theta}_1$ is a much better estimator, as it uses all the data, meaning it will fluctuate much less from sample to sample. $\hat{\theta}_2$ will wildly jump around depending on what that first data point happens to be. ($\mathbb{V}[hat{\theta}_1] = \frac{\mathbb{V}[hat{\theta}_2]}{n}$).

\medskip
\noindent The standard way to measure the quality of estimators by calculating the Expected Squared Error: $\mathbb{E}[||\hat{\theta} - \theta_*||^2]$. In other words, on average, how far away is our estimator's guess from the actual true parameter?
We can break this ESE into two separate parts:
$$\mathbb{E}[||\hat{\theta} - \theta_*||^2] = \mathbb{E}[||\hat{\theta} - \mathbb{E}[\hat{\theta}]||^2] + ||\mathbb{E}[\hat{\theta}] - \theta_*||^2$$
Notice that this translates to: Total Error = Variance + $(\text{Bias})^2$.
\begin{enumerate}
    \item Variance $\mathbb{V}[\hat{\theta}]$: The first term ($\mathbb{E}[||\hat{\theta} - \mathbb{E}[\hat{\theta}]||^2]$) measures how much the estimator's predictions bounce around its own average. It represents the estimator's sensitivity to the specific random sample of data you pulled.
    \item Bias Squared $||\mathbb{E}[\hat{\theta}] - \theta_*||^2$: measures the difference between the estimator's average prediction and the actual true parameter. It tells you if your estimator is systematically missing the mark.
\end{enumerate}
\textbf{What this means for statistical inference}: If an estimator is completely unbiased, its Bias term is exactly 0. Therefore, its total Expected Squared Error is equal only to its Variance ($\mathbb{E}[||\hat{\theta} - \theta_*||^2] = \mathbb{V}[\hat{\theta}]$).\newline
An Efficient Estimator is thus one that is unbiased AND has the minimum variance.

\subsubsection{Fisher Information and Cramér-Rao Lower Bound}
To find an efficient estimator, we need a mathematical floor for the absolute minimum variance.\newline
Firstly, the Fisher Information, denoted as $\mathcal{I}(\theta_*)$ quantifies the amount of information that an observable random variable $X$ carries about an unknown parameter $\theta_*$.\newline
If changing the true parameter drastically changes the probability of our data, the data carries a lot of Fisher information.\newline
It is defined as the expected value of the squared derivative of the log-likelihood:
$$\mathcal{I}(\theta_*) = \mathbb{E} \left[ \left( \frac{\partial}{\partial \theta} \log p_\theta(X) \right)^2 \right]$$
Next, the Cramér-Rao Inequality states that the variance of any unbiased estimator $\hat{\theta}$ cannot possibly be smaller than the inverse of the total Fisher Information provided by the $n$ data points.
$$\mathbb{V}[\hat{\theta}] \ge \frac{1}{n \cdot \mathcal{I}(\theta_*)}$$
Note that as $n \to \infty$ or the data is highly informative ($\mathcal{I}(\theta_*)$ is large), the Cramér-Rao lower bound gets close to zero, meaning it becomes mathematically possible to create a highly precise estimator with very little variance.

\subsection{Mean Estimation}
Given independent observations $X_1, \ldots, X_n \sim P$, find the unknown, true mean of the distribution, $\mu_1 = \mathbb{E}[X_1]$.\newline
Define an estimator, the sample mean
\[\hat{\mu}_1 = \frac{1}{n} \sum_{i=1}^n X_i\]
The sample mean is itself a random variable.\newline
Note that $\mu_1$ refers to the first moment of the distribution, and $X_1, \ldots, X_n$ is i.i.d.\newline
Three properties of a good estimator:
\begin{enumerate}
    \item \textbf{Unbiased}: $\mathbb{E}[\hat{\mu}_{1}]=\mu_{1}$.
    \item \textbf{Variance shrinks with data}: $\mathbb{V}[\hat{\mu}_{1}]=\frac{\mathbb{V}[X_{1}]}{n}$. Shows Effictivity and proves Consistency.
    \item \textbf{Consistent}: $\forall\epsilon>0$, $\lim_{n\rightarrow\infty}P(|\hat{\mu}_{1}-\mu_{1}|\le\epsilon)=1$.
\end{enumerate}
\paragraph{Example} For some distributions, their parameter is equal to the theoretical mean.
\begin{itemize}
    \item Bernoulli distribution $X \sim  Ber(p): \mathbb{E}[X] = p$.
    \item Poisson distribution $X \sim Po(\lambda): \mathbb{E}[X] = \lambda$.
\end{itemize}
Since we know the sample mean is a great estimator for the theoretical mean, we just find the sample mean to estimate the parameters.

\subsection{Variance Estimation}
Given independent observations $X_1, \ldots, X_n \sim P$, estimate the variance $\sigma^2 = \mathbb{E}[(X_1 - \mu_1)^2]$.\newline
Define an estimator, the sample variance
\[\hat{\sigma}_{1}^{2}=\frac{1}{n}\sum_{i=1}^{n}(X_{i}-\hat{\mu}_{1})^{2}\]
This estimator is biased, i.e. $\mathbb{E}[\hat{\sigma}_1^2] = \frac{n-1}{n} \sigma_1^2$.\newline
\textit{Proof.}
\begin{align*}
\mathbb{E}[\hat{\sigma}_1^2] &= \mathbb{E}\left[ \frac{1}{n} \sum_{i=1}^n (X_i - \hat{\mu}_1)^2 \right] \\
&= \frac{1}{n} \mathbb{E}\left[ \sum_{i=1}^n ((X_i - \mu_1) - (\hat{\mu}_1 - \mu_1))^2 \right] && \text{[Inject true mean } \mu_1 \text{]} \\
&= \frac{1}{n} \mathbb{E}\left[ \sum_{i=1}^n (X_i - \mu_1)^2 - 2(\hat{\mu}_1 - \mu_1)\sum_{i=1}^n(X_i - \mu_1) + n(\hat{\mu}_1 - \mu_1)^2 \right] && \text{[Expand quadratic]} \\
&= \frac{1}{n} \mathbb{E}\left[ \sum_{i=1}^n (X_i - \mu_1)^2 - 2n(\hat{\mu}_1 - \mu_1)^2 + n(\hat{\mu}_1 - \mu_1)^2 \right] && \left[\sum(X_i - \mu_1) = n(\hat{\mu}_1 - \mu_1)\right] \\
&= \frac{1}{n} \mathbb{E}\left[ \sum_{i=1}^n (X_i - \mu_1)^2 - n(\hat{\mu}_1 - \mu_1)^2 \right] \\
&= \frac{1}{n} \left[ \sum_{i=1}^n \mathbb{E}[(X_i - \mu_1)^2] - n\mathbb{E}[(\hat{\mu}_1 - \mu_1)^2] \right] && \text{(Linearity of expectation)} \\
&= \frac{1}{n} \left[ n\sigma_1^2 - n\mathbb{V}[\hat{\mu}_1] \right] && \text{(Definitions of } \sigma^2 \text{ and } \mathbb{V}[\hat{\mu}_1]\text{)} \\
&= \frac{1}{n} \left[ n\sigma_1^2 - n\left(\frac{\sigma_1^2}{n}\right) \right] && \text{(Substitute } \mathbb{V}[\hat{\mu}_1] = \frac{\sigma_1^2}{n}\text{)} \\
&= \frac{1}{n} [n\sigma_1^2 - \sigma_1^2] \\
&= \frac{n-1}{n}\sigma_1^2
\end{align*}
Define a corrected estimator:
\[\hat{\tau}^{2}=\frac{n}{n-1}\hat{\sigma}_{1}^{2}=\frac{1}{n-1}\sum_{i=1}^{n}(X_{i}-\hat{\mu}_{1})^{2}\]
\[\mathbb{E}[\hat{\tau}^{2}]= \frac{1}{n-1}(n-1)\sigma_1^2 = \sigma_1^2\]
\textbf{Consistency}: Both $\hat{\sigma}_1^2$ and $\hat{\tau}^2$ are consistent. As $n \to \infty$, the difference between dividing by $n$ versus $n-1$ becomes negligible, and both converge to the true variance.

\subsection{Probability Density Estimation}
Given training data/observations $D = \{X_i\}_{i=1}^n$, $X_i \sim P$, we choose a probabilistic model (hypothesis class of PDFs) $\mathcal{F} = \{p_\theta\}_{\theta \in \Theta}$.\newline
We want to find the specific parameters $\theta$ that make the model $p_\theta$ the best possible approximation of the true distribution $P$ by fitting it to the training data.\newline
For example, a common template to use is the Gaussian distribution. The parameters to tune are the mean and the variance, so as to fit the data.\newline
Another example, Generative AI. Training the model entails analyzing existing examples to find the underlying probability distribution that generated them.
Then during generation phase, the model draws new data samples from those estimated distributions.
This random sample becomes a brand-new, synthetic image or piece of text that looks authentic because it follows the probability rules of the original data.
\begin{figure}[H]
    \centering
    \includegraphics[width=0.8\linewidth]{genai-pdf.png}
    \caption*{The model samples from the estimated distribution.}
    \label{fig:my_image}
\end{figure}

\subsubsection{Maximum Likelihood Estimation (MLE)}
We want to find the parameter $\theta$ that makes the occurrence of $\{X_1, \dots, X_n\}$ as likely as possible, i.e.
\[\max_{\theta} p_\theta(X_1, \dots, X_n)\]
Since we assume the observations are i.i.d., 
\[p_\theta(X_1, \dots, X_n) = \prod_{i=1}^{n} p_\theta(X_i)\]
If $n$ is huge, this number is very small and computers cannot handle it. Hence we take logarithm (recall the property $\log(A \times B) = \log A + \log B$).\newline
Hence, MLE is solving for 
\[\max_{\theta\in\Theta}\sum_{i=1}^{n}\log p_{\theta}(X_{i})\]
For very complex distributions (like deep neural networks), solving this maximization problem directly is often impossible (intractable).\newline
Hence, use approximation methods like EM algorithm, variational Bayes / variational autoencoders (VAEs), diffusion models.

\subsubsection{Kullback–Leibler (KL) divergence}
KL divergence is a statistical measure that quantifies how one probability distribution differs from a second, reference probability distribution.\newline
In estimation context, KL divergence shows how closely our estimated model $q$ matches the true data distribution $p$.\newline
For continuous probability densities $p$ and $q$, the KL divergence between them is
\[D_{KL}(p||q) = \int p(x) \log \frac{p(x)}{q(x)} dx = \mathbb{E}_{X \sim p} \left[ \log \frac{p(X)}{q(X)} \right]\]
Two properties of KL divergence:
\begin{enumerate}
    \item Non-negativity: $D_{KL}(p||q) \ge 0$.
    \item Identity: $D_{KL}(p||q) = 0$ if and only if the two distributions are identical ($p = q$).
\end{enumerate}

\subsubsection{Connection between MLE \& KL divergence}
From the standard MLE sum, multiply by $\frac{1}{n}$ to get the empirical average:
$$\max_{\theta} \frac{1}{n} \sum_{i=1}^{n} \log p_{\theta}(X_{i})$$
As $n \to \infty$, the empirical average converges to the theoretical expected value over the true data distribution $p$, by Law of Large Numbers,
$$\max_{\theta} \mathbb{E}_{X \sim p}[\log p_{\theta}(X)]$$
To connect this to KL divergence, first subtract a constant.
\begin{align*}
    \mathbb{E}_{X \sim p}[\log p_{\theta}(X)] &\sim \mathbb{E}_{X \sim p}[\log p_{\theta}(X)] - \mathbb{E}_{X \sim p}[\log p(X)] \\
    &= \mathbb{E}_{X \sim p}\left[ \log \frac{p_{\theta}(X)}{p(X)} \right] \\
    &= -D_{KL}(p || p_{\theta})
\end{align*}
Maximising the log-likehood minimises the KL divergence.

\subsubsection{Example - Bernoulli Distribution}
Given random variable $X \sim Ber(\theta)$, $P_\theta(X = 1) = \theta$, $P\theta(X = 0) = 1 - \theta$,
\[P_\theta(X) = \theta^X \cdot (1 - \theta)^{1 - X}\]
Observations are $X_i \sim P$, ($i = 1, 2, \ldots, n$), which are i.i.d.
The likelihood function is
\[L(\theta) = \prod_{i=1}^{n} \theta^{X_i} (1-\theta)^{1-X_i}\]
Since the product is hard to calculate, the log likelihood function is
\[\mathcal{L}(\theta) = \sum_{i=1}^{n} \left[ X_i \log(\theta) + (1-X_i) \log(1-\theta) \right]\]
To find $\theta$ that maximises the function, take the derivative wrt $\theta$ and set to 0:
\[\frac{d\mathcal{L}}{d\theta} = \frac{1}{\theta} \sum_{i=1}^{n} X_i - \frac{1}{1-\theta} \sum_{i=1}^{n} (1-X_i) = 0\]
Solving, we get
\[\hat{\theta} = \frac{1}{n} \sum_{i=1}^{n} X_i\]
The most optimal way to estimate the probability here is just the sample mean.\newline
For example, given 10 coin flips, $D = \{1, 1, 0, 1, 1, 1, 0, 1, 1, 0\}$,
$\hat{\theta} = 0.7$.\newline
Interpretation: The parameter that makes this specific dataset the most mathematically likely to occur is $\theta = 0.7$.

\smallskip
\noindent \textbf{Is this estimator efficient? (has the lowest possible variance)}\newline
For a Bernoulli distribution, the variance of a single trial $X_i$ is known to be $\mathbb{V}[X_i] = \theta_*(1-\theta_*)$.
Then,
\begin{align*}
    \mathbb{V}[\hat{\theta}] &= \mathbb{V}\left[\frac{1}{n}\sum_{i=1}^n X_i\right] \\
    &= \frac{1}{n^2} \sum_{i=1}^n \mathbb{V}[X_i] \\
    &= \frac{\theta_*(1-\theta_*)}{n}
\end{align*}
We start with the log-likelihood of a single Bernoulli trial:
$$\log P(X|\theta_*) = X\log(\theta_*) + (1-X)\log(1-\theta_*)$$
Using the definition of Fisher information:
\begin{align*}
    \mathcal{I}(\theta_*) &= \mathbb{E}\left[\left(\frac{\partial}{\partial \theta_*} \log P(X|\theta_*)\right)^2\right] \\
    &= \mathbb{E}\left[\left(\frac{X}{\theta_*} - \frac{1-X}{1-\theta_*}\right)^2\right] \\
    &= \mathbb{E}\left[\frac{X^2}{\theta_*^2} - 2\frac{X(1-X)}{\theta_*(1-\theta_*)} + \frac{(1-X)^2}{(1-\theta_*)^2}\right] \\
    &= \mathbb{E}\left[ \frac{X}{\theta_*^2} - 0 + \frac{1-X}{(1-\theta_*)^2} \right] \\
    &= \frac{\mathbb{E}[X]}{\theta_*^2} + \frac{\mathbb{E}[1-X]}{(1-\theta_*)^2} \\
    &= \frac{\theta_*}{\theta_*^2} + \frac{1-\theta_*}{(1-\theta_*)^2} \\
    &= \frac{1}{\theta_*(1-\theta_*)}
\end{align*}
Then the Cramér-Rao Lower Bound is:
\begin{align*}
    \text{CRLB} &= \frac{1}{n \cdot \mathcal{I}(\theta_*)} \\
    &= \frac{1}{n \cdot \left(\frac{1}{\theta_*(1-\theta_*)}\right)} \\
    &= \frac{\theta_*(1-\theta_*)}{n} \\
    &= \mathbb{V}[\hat{\theta}]
\end{align*}
\textbf{The implication}:  $\hat{\theta} = \frac{1}{n}\sum X_i$ is an efficient estimator.

\subsubsection{Example - Multivariate Gaussian Distribution}
Moving from just a single variable $x$, we now consider data points in $d$-dimensional space $\mathbb{R}^d$.\newline
Define the mean vector $\mu$, which is a vector that contains the average value for every single dimension, and $\mu = \mathbb{E}_{X \sim p_\theta}[X]$.\newline
Also, define the covariance matrix $\Sigma \in \mathbb{R}^{d \times d}$, which is a square matrix that captures the variance. Here, $\Sigma = \mathbb{E}_{X \sim p_\theta}[(X - \mu)(X - \mu)^\top]$.
It tells us the variance of each individual dimension and how all the dimensions correlate with one another.\newline
Recall the $d$-variate Gaussian PDF:
\[p_\theta(x) = \frac{1}{(2\pi)^{d/2}\sqrt{\det \Sigma}} \exp\left(-\frac{1}{2}(x - \mu)^\top \Sigma^{-1} (x - \mu)\right)\]
with the normalisation constant in the front, and the term in the exponent calculating the "distance" of point $x$ from the mean $\mu$, scaled by the covariance.\newline
The goal is to use MLE to calculate $\mu$ and $\Sigma$ that best fit the data.\newline
\begin{align*}
\mathcal{L}(\mu, \Sigma) &= \log \prod_{i=1}^{n} \frac{1}{(2\pi)^{d/2} |\Sigma|^{1/2}} \exp\left(-\frac{1}{2}(X_i - \mu)^\top \Sigma^{-1} (X_i - \mu)\right) && \text{(Log of the joint PDF)} \\
&= \sum_{i=1}^{n} \left[ -\frac{d}{2}\log(2\pi) - \frac{1}{2}\log|\Sigma| - \frac{1}{2}(X_i - \mu)^\top \Sigma^{-1} (X_i - \mu) \right] && \text{(Log rules: product to sum)} \\
&= -\frac{nd}{2}\log(2\pi) - \frac{n}{2}\log|\Sigma| - \frac{1}{2}\sum_{i=1}^{n} (X_i - \mu)^\top \Sigma^{-1} (X_i - \mu) && \text{(Factor constants out of sum)}
\end{align*}
$$(\hat{\mu}, \hat{\Sigma}) = \arg\max_{\mu \in \mathbb{R}^d, \Sigma \in \mathbb{R}^{d \times d}} -\frac{n}{2}\log(\det\Sigma) - \frac{1}{2}\sum_{i=1}^{n}(X_i - \mu)^\top\Sigma^{-1}(X_i - \mu)$$
Taking derivative of this function w.r.t. to $\mu$ and $\Sigma$ and setting it to zero,
\begin{itemize}
    \item $\hat{\mu} = \frac{1}{n}\sum_{i=1}^{n}X_i$
    \item $\hat{\Sigma} = \frac{1}{n}\sum_{i=1}^{n}(X_i - \hat{\mu})(X_i - \hat{\mu})^\top$
\end{itemize}
\textbf{Conclusion}: Applying MLE to a multivariate Gaussian distribution gives us the sample mean and the sample covariance matrix as our optimal estimators.
As the number of samples increases, these calculated estimates will converge to the true, underlying parameters $\mu$ and $\Sigma$ of the distribution.

\subsection{Bayesian Inference}
Methods like MLE treated the parameter $\theta$ as a fixed, absolute, but unknown constant to solve for.\newline
Bayesian Inference instead considers a probability distribution of $\theta$.
$\theta$ is treated a random variable with its own probability distribution, and we can update our entire belief about the parameter as we gather more data.\newline
The probabilistic structure is:
\begin{itemize}
    \item Prior Distribution, $p(\theta)$: initial belief about the distribution of the parameter before you have observed any data.
    \item Likelihood, $p(x|\theta)$: how likely you are to observe the data $x$, assuming a specific parameter $\theta$ is true.
    \item Joint Distribution, $p(x, \theta) = p(x|\theta)p(\theta)$: combined probability of both the parameter and the data occurring.
\end{itemize}
The goal is to infer $\theta$ by calcuting the posterior distribution using Bayes' Theorem:
$$p(\theta|X_1, ..., X_n) = \frac{p(X_1, ..., X_n|\theta)p(\theta)}{p(X_1, ..., X_n)} = \frac{\prod_i p(X_i|\theta)p(\theta)}{\int \prod_i p(X_i|\theta')p(\theta')d\theta'}$$
The numerator can be written as the product of individual likelihoods since the observations are independent. The $\theta$ here represents the specific hypothesis we are currently testing.\newline
The denominator calculates the total marginal probability of observing this specific set of data. We use the dummy variable $\theta'$ here to show that this integral sweeps across every possible value of the parameter to calculate a grand total. Because it integrates out the parameter, the denominator evaluates to a flat constant number.\newline
With this updated belief, we want to make a prediction about new, unseen data. In other words, given all the past data observations $X_1, ..., X_n$ I have already seen, what is the probability of observing a completely new data point $x$?\newline
We want to solve for $\hat{p}(x)$, the posterior predictive distribution.
$$\hat{p}(x) = \mathbb{E}_{\theta\sim p(\theta|X_1, ..., X_n)}[p(x|\theta)]$$
The expected value is calculating the weighted average of the likelihood $p(x|\theta)$, where the "weights" are determined by how plausible each $\theta$ is according to your posterior distribution $p(\theta|X_1, ..., X_n)$.\newline
Recall that if we used MLE, we would have calculated the single best parameter $\hat{\theta}_{MLE}$, and to predict a new data point, we just plug it into $p(x|\hat{\theta}_{MLE})$.\newline

\medskip
\noindent The expected value of a continuous function is found by integrating it over the probability distribution. Multiplying ($p(x|\theta)$) by the weights ($p(\theta|X_1, ..., X_n)$) and integrating over $\theta$,
\begin{align*}
    \hat{p}(x) &= \int p(x|\theta) \cdot p(\theta|X_1, ..., X_n) d\theta \\
    &= \int p(x|\theta) \left[ \frac{\prod_i p(X_i|\theta)p(\theta)}{\int \prod_i p(X_i|\theta')p(\theta')d\theta'} \right] d\theta \\ 
    &= \frac{\int p(x|\theta) \prod_i p(X_i|\theta)p(\theta)d\theta}{\int \prod_i p(X_i|\theta')p(\theta')d\theta'}
\end{align*}
\textbf{Conclusion}: Because Bayesian inference averages over all possible parameters rather than relying on a single estimate, its predictions are much safer, especially when you have very little data.
If your dataset is small, your posterior distribution will be wide (highly uncertain). The integral naturally accounts for this uncertainty, preventing you from making overconfident predictions based on a small sample size.
As you gather more data ($n \rightarrow \infty$), your posterior distribution becomes an incredibly narrow spike around the true parameter, and this massive integral mathematically simplifies down to become virtually identical to the MLE approach.

\subsubsection{Bayesian Inference with Gaussian Distribution}
Calculating that massive integral in the denominator to find the posterior is often computationally impossible for complex distributions.\newline
Instead, we can the use a property of Gaussians: if your prior is a Gaussian, and your likelihood is a Gaussian, your resulting posterior is guaranteed to also be a Gaussian.\newline
First, model the prior and likelihood by Gaussian distributions:
\begin{itemize}
    \item The Prior, $\theta \sim \mathcal{N}(\mu_0, \sigma_0^2)$: Before seeing any data, you believe the true parameter $\theta$ is somewhere around $\mu_0$, but you have some uncertainty, represented by the variance $\sigma_0^2$.
    \item The Likelihood ($X_i \sim \mathcal{N}(\theta, \sigma^2)$): Observing data points $X_1, ..., X_n$, under the assumption that these data points are generated by a Gaussian distribution centered around the true parameter $\theta$, with a known variance of $\sigma^2$.
\end{itemize}
After observing $n$ data points, the belief updates.
The posterior variance $\hat{\sigma}_n^2$ and posterior mean $\hat{\mu}_n$ are
$$\hat{\sigma}_n^2 = \frac{\sigma^2 \sigma_0^2}{n\sigma_0^2 + \sigma^2} \quad \hat{\mu}_n = \frac{\sigma^2\sigma_0^2}{n\sigma_0^2 + \sigma^2} \left( \frac{\sum_{i=1}^{n}X_i}{\sigma^2} + \frac{\mu_0}{\sigma_0^2} \right)$$
As $n \to \infty$,
$$\lim_{n \rightarrow \infty} \hat{\sigma}_n^2 = 0 \quad \lim_{n \rightarrow \infty} \hat{\mu}_n = \frac{\sum_{i=1}^{n}X_i}{n}$$
The uncertainty drops to 0, and the mean becomes exactly equal to the sample mean, ignoring the initial belief.\newline

\medskip
\noindent\textbf{Conclusion}: Recall that the sample mean $\frac{\sum X_i}{n}$ is the exact solution you get when you use MLE.\newline
Bayesian Inference is asymptotically equivalent to MLE.\newline
In small data regime (we have little data), MLE is vulnerable to extreme outliers and overfitting. Bayesian inference instead anchors the prediction to your prior belief.\newline
In large data regime, the sheer volume of evidence overwhelms whatever initial bias or belief you had. The prior's influence vanishes, and the Bayesian approach perfectly mathematically aligns with the Frequentist (MLE) approach.

\subsubsection{Maximum a Posteriori Estimation}
In Bayesian inference, the perfect way to make a prediction is to calculate $\hat{p}(x) = \mathbb{E}_{\theta \sim p(\theta|X_1, ..., X_n)}[p(x|\theta)]$.\newline
However, this is intractable. This means that for complex, real-world models (like deep neural networks), the integrals required to calculate the full posterior and the expected value are too complicated and computationally heavy for a computer to actually solve.\newline
Instead of averaging all possible parameters, MAP Estimation simply finds the single most likely parameter value at the very peak of the posterior distribution curve.\newline
This value is denoted as $\hat{\theta}_{MAP}$, and it is found by taking the $\arg\max$ of the posterior:
$$\hat{\theta}_{MAP} = \arg\max_\theta p(\theta|X_1, ..., X_n)$$
Specifically, the core MAP equation is
$$\max_{\theta \in \Theta} \left\{ \sum_{i=1}^n \log p(X_i|\theta) + \log p(\theta) \right\}$$
There are two distinct parts in this equation:
\begin{enumerate}
    \item The data likelihood $\sum_{i=1}^n \log p(X_i|\theta)$: This is the same log-likelihood term we maximized in standard Maximum Likelihood Estimation (MLE). It pushes the model to pick parameters that fit the observed data well.
    \item The prior as a penalisation term $\log p(\theta)$: This is the log of your prior distribution acting as a penalisation term. If the MLE part tries to pick an extreme parameter that perfectly fits the data but goes against your initial beliefs, this prior term will mathematically penalize (reduce) the overall score.
\end{enumerate}
Because of this added prior term, this problem is also called penalized maximum likelihood estimation.\newline
Also, as $n \to \infty$, the data likelihood term (which is a sum) becomes massive, completely overwhelming the single, static $\log p(\theta)$ prior term. With enough data, MAP becomes identical to MLE.\newline
\newline 
After finding this single best parameter, we can use it to make predictions as $p(x|\hat{\theta}_{MAP})$.\newline
MAP uses Bayesian logic (incorporating a prior belief) but outputs a single parameter like the Frequentist MLE approach.




\end{document}